\textbf{Background:} Standard depth-damage functions used in climate risk assessment---including those underpinning CLIMADA and JRC impact models---treat flood damage as an instantaneous function of water depth, ignoring the duration of inundation.

\textbf{Methods:} We extend the generalised logistic (Richards) damage function with duration-dependent parameters and calibrate it on 295{,}637 flood insurance claims from the FEMA National Flood Insurance Program (OpenFEMA snapshot 2026-05-09), spanning three major US hurricanes: Sandy (2012, $\sim$36\,h), Harvey (2017, $\sim$288\,h), and Katrina (2005, $\sim$120\,h). To bridge the residential calibration to non-residential assets, we introduce a sector adjustment multiplier $\kappa$ calibrated on 82{,}650 NFIP non-residential claims grouped by building stories and insurance coverage, combined with a coverage-gradient extrapolation to corporate scale.

\textbf{Results:} Prolonged flooding increases structural damage by a factor of 2.6$\times$ at the same depth, and this effect operates primarily through a shift in the curve's inflection point ($\lambda_a = 0.49$), while the fundamental shape of the damage function remains universal across events. The calibrated model reduces mean absolute error by about 50\% relative to JRC depth-damage curves and achieves Spearman $\rho = 0.949$ ($p < 0.001$) against actual financial damage across 301 ZIP codes. Out-of-sample validation on 24 companies across six flood events yields a median predicted-to-actual ratio of 0.90$\times$ (23/24 within 0.4--2.5$\times$), with no corporate loss disclosures entering the calibration.

\textbf{Conclusions:} Duration is a first-order determinant of flood damage that existing vulnerability frameworks systematically miss. We demonstrate the full pipeline on Archer Daniels Midland's 12 flood-exposed Midwest facilities, producing facility-level loss estimates with duration effects, CapEx/OpEx decomposition, and insurance offsets. The framework translates physical damage into shifts in probability of default via an interest coverage ratio approach, enabling direct integration with banks' internal credit risk models.
